\chapter{The master's thesis: a defensible contribution}
\label{ch:thesis}

\section{A motivating problem}
A thesis proposal says, \enquote{We will build an innovative AI platform that transforms public services}. It contains ambition but no feasible boundary, research question, method, contribution, or evaluation. A master's thesis earns its strength through delimitation: a claim small enough to investigate and important enough to matter.

\learningobjectives{You should be able to select and delimit a feasible thesis problem; build a conceptual framing and literature review; choose a research or design method; plan an artefact and evaluation; distinguish contribution types; report results and limitations; manage scholarly collaboration; and present a defensible argument.}

\section{The five-minute explanation}
A thesis is a sustained argument. It tells the reader what problem was investigated, why it matters, what is already known, what was done, what was found, what can be concluded, and what remains uncertain. A software artefact may be central, but the thesis must explain what knowledge the work contributes and how evidence supports that claim.

\section{From topic to research question}
Start with a domain problem and an accessible source of evidence. Delimit by population, context, task, system boundary, time, data, and contribution. A feasible question often has the form:
\begin{quote}
How does [intervention or design] affect or support [construct or activity] for [population] in [context], and under which conditions?
\end{quote}
The wording should match the design. If there is no comparison, do not imply a causal effect. If the study concerns one organisation, delimit transfer.

\section{Literature and conceptual framing}
The literature review should justify definitions, mechanisms, methods, and gap. A gap may be empirical, theoretical, methodological, design-oriented, or contextual. It need not mean that nobody has ever studied the topic. The contribution can be a careful replication, synthesis, boundary condition, or evaluated artefact.

Build a matrix with source, question, theory, method, data, finding, limitation, and relevance. Cite primary authors for concepts and original empirical studies where possible.

\section{Method and artefact}
Possible thesis designs include:
\begin{itemize}
\item an empirical study of use, adoption, or organisational change;
\item a design-science project that builds and evaluates an artefact;
\item a theoretical investigation of an interactive-system problem;
\item a mixed-methods study integrating qualitative context and quantitative outcomes;
\item a systematic or scoping review with a transparent protocol.
\end{itemize}
An artefact can be a software system, interaction design, model, method, dataset, or architecture. Specify the build process, requirements, design rationale, implementation limits, evaluation, and reproducibility.

\section{Contribution types}
\begin{longtable}{p{0.22\textwidth}p{0.32\textwidth}p{0.34\textwidth}}
\caption{Distinguishing thesis contribution types.}\\
\toprule
\textbf{Contribution} & \textbf{What it adds} & \textbf{Evidence required}\\
\midrule
Empirical & A finding about a defined context or population & Data, design, analysis, uncertainty, limitations\\
Theoretical & A concept, relationship, mechanism, or boundary condition & Conceptual argument and engagement with literature\\
Methodological & A procedure or measurement approach & Rationale, implementation, comparison or evaluation\\
Artefact/design & A system, prototype, architecture, or design knowledge & Requirements, construction, rationale, evaluation\\
Replication & A reasoned reproduction or extension of prior work & Faithful protocol, comparison, deviations\\
Synthesis & An integrated account of prior evidence & Search, selection, appraisal, and synthesis method\\
\bottomrule
\end{longtable}

A thesis can contain several contribution types, but name which is central. Novelty is not a substitute for validity.

\section{Results, discussion, and limitations}
Results report what the analysis produced. Discussion interprets the result in relation to the question and literature. Limitations explain how the design constrains inference. A limitation should not be a ritual list; it should say which claim is weakened and in what direction.

A defensible conclusion uses calibrated verbs: \emph{describes}, \emph{is consistent with}, \emph{suggests}, \emph{supports under these conditions}, or \emph{does not establish}. Avoid \emph{proves} unless the domain and proof standard genuinely warrant it.

\section{Planning and supervision}
Create a reverse plan from the final delivery and presentation. Reserve time for ethics, access, recruitment, data cleaning, implementation debt, analysis, writing, references, visual QA, and contingency. Supervision is a scholarly collaboration, not an outsourced project-management function. Bring concrete questions, evidence of progress, and a record of decisions.

\section{Writing and presenting the thesis}
A clear chapter structure is an argument structure: introduction, related work, method, design or system, results, discussion, conclusion, references, appendices. Figures should be introduced, captioned, interpreted, and readable in print. Code should be versioned and runnable or its irreproducibility explained.

When presenting a thesis, practise answers to:
\begin{enumerate}
\item What is the exact research question?
\item What is the central contribution?
\item Why is the method appropriate?
\item What alternative explanation is most serious?
\item Which design decision would you change with more time?
\item What does the thesis not establish?
\end{enumerate}

\section{Project application}
The thesis should connect the full progression of this book: problem, model, requirements, design, implementation, data, evaluation, organisation, and claim. Its scope, contribution, method, and evidence should be clear enough that a critical reader can reconstruct the path from question to conclusion.

\section{Summary and key terms}
A master's thesis is a bounded argument supported by theory, method, artefact or data, analysis, and reflection. Delimitation is a strength. Contribution types differ in evidence requirements. Results, interpretation, and limitation should be separate. A strong presentation demonstrates understanding and judgment, not only a polished final product.

\begin{keyterms}thesis; delimitation; research question; literature review; conceptual framework; design science; artefact; empirical contribution; theoretical contribution; methodological contribution; replication; synthesis; result; discussion; limitation; reproducibility; presentation.\end{keyterms}

\section{Exercises and worked solutions}
\exercise{Why is \enquote{build an innovative AI platform} not yet a thesis question?}
\solution{It names an artefact and an evaluative adjective but not a phenomenon, population, context, comparison, method, or contribution. A bounded version might ask how a specified AI interaction affects a defined task for a defined group under specified conditions, with an explicit evaluation design.}

\exercise{What is the difference between a limitation and a failure?}
\solution{A limitation is a condition that restricts what can be inferred, even when the study is competently executed. A failure is a departure from the planned or required process that may invalidate part of the evidence. A limitation can be responsibly designed; a failure must be analysed and disclosed.}

\exercise{What should a supervisor meeting produce?}
\solution{It should produce clarified questions, decisions, risks, next actions, and a record of unresolved issues. The student remains responsible for the argument and should not treat disagreement as a substitute for methodological reasoning.}

\section{Further reading}
Use the research-methods chapter as the methodological foundation. For case-study reporting, see \citet{runeson2009case}; for experimentation, see \citet{wohlin2012experimentation}.
